% Section 4: Axiom Verification
% \input-able LaTeX file for Paper 5
% Conventions: a \sp b = sequential product (dot notation matching vdW)
%   C_p = Alfsen-Shultz compression (P-projection)
%   P_{ij} = Peirce 1-space projection
%   All axiom statements from arXiv:1803.11139 Definition 2 EXCLUSIVELY

\section{Axiom Verification}\label{sec:axioms}

We now verify that the self-modeling sequential product defined in
Section~\ref{sec:sp} satisfies all seven axioms of van de
Wetering's Definition~2~\cite{vandeWetering2019b}. By Theorem~1 of~\cite{vandeWetering2019b},
this establishes that the state space carries the structure of a Euclidean
Jordan algebra.

Throughout this section, let $V$ be a finite-dimensional spectral order
unit space and let the sequential product be given by the corrected
formula~\eqref{eq:corrected-product}:
\begin{equation}\label{eq:sp-explicit}
  \sp{a}{b}
  = \sum_i \lambda_i\, C_{p_i}(b)
  + \sum_{i<j} \sqrt{\lambda_i \lambda_j}\, P_{ij}(b),
\end{equation}
where $a = \sum_i \lambda_i\, p_i$ is the spectral decomposition of $a$,
$C_{p_i}$ is the Alfsen--Shultz compression for the face generated by
$p_i$, and $P_{ij}(b) = b - \sum_k C_{p_k}(b)$ is the Peirce $1$-space
projection.

%--------------------------------------------------------------------
\subsection*{S1 (Additivity)}

\begin{quote}
\emph{The map $b \mapsto \sp{a}{b}$ is additive:
$\sp{a}{(b+c)} = \sp{a}{b} + \sp{a}{c}$
whenever $b + c \le \mathbf{1}$.}
\end{quote}

Each compression $C_{p_i}$ is a positive linear map on $V$
(Alfsen--Shultz~\cite{AlfsenShultz2003}, Ch.~7), and the Peirce
projection $P_{ij} = \mathrm{id} - \sum_k C_{p_k}$ is a difference of
linear maps, hence linear. Since the coefficients $\lambda_i$ and
$\sqrt{\lambda_i\lambda_j}$ depend only on $a$, the map $b \mapsto
\sp{a}{b}$ is a fixed linear combination of linear maps applied to $b$,
and therefore additive. \qed

%--------------------------------------------------------------------
\subsection*{S2 (Continuity)}

\begin{quote}
\emph{The map $a \mapsto \sp{a}{b}$ is continuous in the order unit
norm.}
\end{quote}

In a finite-dimensional spectral order unit space, the spectral
decomposition $a = \sum_i \lambda_i(a)\, p_i(a)$ varies continuously with
$a$ (Alfsen--Shultz~\cite{AlfsenShultz2003}, Ch.~9; in finite dimensions,
eigenvalues are continuous functions and the spectral functional calculus
extends continuously through degeneracies). The scalar maps
$\lambda_i \mapsto \lambda_i$ and
$(\lambda_i,\lambda_j)\mapsto\sqrt{\lambda_i\lambda_j}$ are continuous on
$[0,1]$. Since the compressions $C_{p_i}$ and Peirce projections $P_{ij}$
depend continuously on the projectors $p_i$, the map $a \mapsto \sp{a}{b}$
is a composition of continuous maps and hence continuous. \qed

%--------------------------------------------------------------------
\subsection*{S3 (Unitality)}

\begin{quote}
\emph{$\sp{\mathbf{1}}{a} = a$ for all effects $a$.}
\end{quote}

The order unit $\mathbf{1}$ has the trivial spectral decomposition
$\mathbf{1} = 1\cdot\mathbf{1}$ (a single eigenvalue with the maximal
projective unit). The Peirce $1$-space sum is empty, and the compression
$C_{\mathbf{1}}$ for the full face is the identity map on $V$
(Alfsen--Shultz~\cite{AlfsenShultz2003}, Def.~7.1). Therefore
$\sp{\mathbf{1}}{a} = 1 \cdot C_{\mathbf{1}}(a) = a$.

This is the axiom whose verification critically depends on the Peirce
$1$-space correction: the naive product $\sum_i \lambda_i\, C_{p_i}(b)$
(without the $\sqrt{\lambda_i\lambda_j}\,P_{ij}(b)$ terms) gives
$\sp{\mathbf{1}}{a} = \mathrm{pinch}(a) \ne a$ for effects with
off-diagonal components, because the sum of compressions over a resolution
of unity is the pinching map, not the identity. The self-modeling feedback
that introduces the Peirce $1$-space terms is what restores unitality. \qed

%--------------------------------------------------------------------
\subsection*{S4 (Orthogonality Symmetry)}\label{subsec:S4}

\begin{quote}
\emph{If $\sp{a}{b} = 0$ then $\sp{b}{a} = 0$.}
\end{quote}

This is the key axiom. The remaining axioms (S1--S3, S5--S7) follow from
linearity and spectral calculus with comparatively short arguments. S4
requires a substantive proof that exploits the facial structure of the
order unit space. We give the argument in full; see
Appendix~\ref{app:S4-proof} for additional details.

\begin{theorem}\label{thm:S4}
Let $(V, \le, \mathbf{1})$ be a finite-dimensional spectral order unit
space equipped with the sequential product~\eqref{eq:sp-explicit}. If\,
$\sp{a}{b} = 0$ for effects $a,b \in [0,\mathbf{1}]_V$, then
$\sp{b}{a} = 0$.
\end{theorem}

\begin{proof}[Proof sketch]
Write $a = \sum_{i=1}^n \lambda_i\, p_i$ (spectral decomposition). Two
cases arise.

\medskip
\noindent\textbf{Case~A (full-rank).}
Suppose all $\lambda_i > 0$. The condition $\sp{a}{b} = 0$ and the direct
sum structure of the Peirce decomposition
\[
  V = \bigoplus_i V_2(p_i) \;\oplus\; \bigoplus_{i<j} V_1(p_i,p_j)
\]
(Alfsen--Shultz~\cite{AlfsenShultz2003}, Theorem~9.37)
force each component to vanish individually: $\lambda_i\, C_{p_i}(b) = 0$
and $\sqrt{\lambda_i\lambda_j}\, P_{ij}(b) = 0$. Since $\lambda_i > 0$
for all~$i$, the completeness of the Peirce decomposition gives $b = 0$,
so $\sp{b}{a} = 0$ trivially.

\medskip
\noindent\textbf{Case~B (rank-deficient).}
Suppose $\lambda_1,\ldots,\lambda_m > 0$ and
$\lambda_{m+1} = \cdots = \lambda_n = 0$ for some $1 \le m < n$. As
before, the Peirce direct sum forces $C_{p_i}(b) = 0$ for each $i \le m$.
The crucial step invokes the Alfsen--Shultz facial absorption theorem
(Proposition~7.43 of~\cite{AlfsenShultz2003}):
\begin{quote}
  \emph{If $C_p(b) = 0$ and $b \ge 0$, then
  $b \in \mathrm{face}(p^\perp)$.}
\end{quote}
Applying this to each $p_i$ ($i \le m$) places $b$ in the complementary
face $\mathrm{face}\!\left(\sum_{k=m+1}^n p_k\right)$. By the facial
structure of spectral order unit spaces, the Peirce $1$-space components
$V_1(p_i, p_j)$ connecting a face to its complement are excluded: an
effect in $\mathrm{face}(p^\perp)$ has zero component in every
$V_1(p_i,p_j)$ for $i\le m$, $j > m$. Therefore $b$ is supported
entirely on the zero-eigenvalue block of $a$.

For the reverse product $\sp{b}{a}$, write
$b = \sum_j \mu_j\, q_j$ with each $q_j$ ($\mu_j > 0$) lying in
$\mathrm{face}\!\left(\sum_{k>m} p_k\right)$. Since $a$ is supported on
the complementary face $\mathrm{face}\!\left(\sum_{i\le m}
p_i\right)$, the facial orthogonality theorem gives $C_{q_j}(a) = 0$ for
all $j$ with $\mu_j > 0$. The Peirce $1$-space terms $Q_{jk}(a)$ either
vanish by facial structure (when both $\mu_j,\mu_k>0$) or carry zero
weight $\sqrt{\mu_j\mu_k}=0$. Hence $\sp{b}{a} = 0$.
\end{proof}

\begin{remark}[Independence of $\varphi$]\label{rem:S4-phi-independent}
Axiom S4 is $\varphi$-independent: it holds for \emph{any} mixing
function $f$ with $f(0,x) = 0$, not just for the faithful choice
$f = \sqrt{\lambda_i\lambda_j}$. The proof depends only on the facial
geometry of the order unit space and the vanishing of $f$ at zero
eigenvalues, not on the specific form of $f$ for positive arguments.
Consequently, S4 is a structural property of spectral order unit spaces,
not an artifact of the self-modeling construction.
\end{remark}

%--------------------------------------------------------------------
\subsection*{S5 (Compatible Associativity)}

\begin{quote}
\emph{If $a$ and $b$ are compatible (i.e., $\sp{a}{b} = \sp{b}{a}$),
then $\sp{a}{(\sp{b}{c})} = \sp{(\sp{a}{b})}{c}$ for all effects $c$.}
\end{quote}

Compatible effects share a joint spectral decomposition. For sharp
compatible effects $p,q$, the compressions commute: $C_p \circ C_q = C_q
\circ C_p$ (Alfsen--Shultz~\cite{AlfsenShultz2003}, Prop.~7.49), and
their composition satisfies $C_p \circ C_q = C_{p \wedge q}$
(Prop.~7.50). For general compatible effects, simultaneous
diagonalizability reduces the product to a function of shared spectral
values, and the identity
$\sqrt{\alpha\beta\cdot\gamma\delta}
= \sqrt{\alpha\gamma}\cdot\sqrt{\beta\delta}$
for non-negative reals ensures compatibility of the mixing function. \qed

%--------------------------------------------------------------------
\subsection*{S6 (Additivity of Compatibility)}

\begin{quote}
\emph{If $a \compatible b$ then $a \compatible (\mathbf{1} - b)$; if
also $a \compatible c$ then $a \compatible (b + c)$ whenever
$b + c \le \mathbf{1}$.}
\end{quote}

\noindent\textbf{Part (i): $a \compatible b \implies a \compatible
(\mathbf{1} - b)$.} By S1, the map $b \mapsto \sp{a}{b}$ is additive,
and by construction $\sp{a}{\mathbf{1}} = a$ (since
$C_{p_i}(\mathbf{1}) = p_i$ and $P_{ij}(\mathbf{1}) = 0$). Therefore
$\sp{a}{(\mathbf{1} - b)} = a - \sp{a}{b}$. For the reverse product:
when $a$ and $b$ are compatible ($[a,b] = 0$ in the Jordan algebra),
the Luders product reduces to the ordinary product, giving
$\sp{b}{a} = ba$ and $\sp{(\mathbf{1} - b)}{a} = (1-b)a = a - ba$.
Since $a \compatible b$ means $\sp{a}{b} = \sp{b}{a}$, we have
$\sp{a}{(\mathbf{1}-b)} = a - \sp{a}{b} = a - \sp{b}{a} =
\sp{(\mathbf{1}-b)}{a}$, establishing $a \compatible (\mathbf{1}-b)$.

\noindent\textbf{Part (ii): $a \compatible b$ and $a \compatible c
\implies a \compatible (b+c)$.} By S1,
$\sp{a}{(b+c)} = \sp{a}{b} + \sp{a}{c}$. For the reverse: since $a$
is compatible with both $b$ and $c$, all three are simultaneously
diagonalizable. In the common eigenbasis the Luders product reduces to
ordinary multiplication, so $\sp{(b+c)}{a} = (b+c)\,a = ba + ca =
\sp{b}{a} + \sp{c}{a}$. Using compatibility:
$\sp{a}{(b+c)} = \sp{a}{b} + \sp{a}{c} = \sp{b}{a} + \sp{c}{a} =
\sp{(b+c)}{a}$. \qed

%--------------------------------------------------------------------
\subsection*{S7 (Multiplicativity of Compatibility)}

\begin{quote}
\emph{If $a \compatible b$ and $a \compatible c$ then
$a \compatible \sp{b}{c}$.}
\end{quote}

By the functional calculus, $[a,b] = 0$ implies $[a, f(b)] = 0$ for any
continuous $f$, and in particular $[a, \sqrt{b}] = 0$. Now
$\sp{b}{c} = \sqrt{b}\,c\,\sqrt{b}$ (Luders form). Since $a$ commutes
with $\sqrt{b}$ and with $c$:
\[
  a\,\sp{b}{c}
  = a\,\sqrt{b}\,c\,\sqrt{b}
  = \sqrt{b}\,a\,c\,\sqrt{b}
  = \sqrt{b}\,c\,a\,\sqrt{b}
  = \sqrt{b}\,c\,\sqrt{b}\,a
  = \sp{b}{c}\,a.
\]
Therefore $[a, \sp{b}{c}] = 0$, and since the Luders product of
commuting effects equals their ordinary product, both
$\sp{a}{(\sp{b}{c})}$ and $\sp{(\sp{b}{c})}{a}$ equal
$a \cdot \sp{b}{c}$, establishing $a \compatible \sp{b}{c}$.

At the OUS level, the argument uses compression commutativity
(Niestegge~\cite{Niestegge2009}, Lemma~3.3): when the compressions for
$a$'s spectral projectors commute with those for $b$'s and $c$'s, they
commute with every building block of $\sp{b}{c}$. \qed

%--------------------------------------------------------------------
\bigskip

With S1--S7 verified, we invoke the central classification result:

\begin{theorem}[{EJA classification; \cite[Theorem~1]{vandeWetering2019b}}]
\label{thm:EJA}
Let $(V, \le, \mathbf{1})$ be a finite-dimensional order unit space
admitting a sequential product satisfying S1--S7. Then $V$ is
order-isomorphic to a Euclidean Jordan algebra.
\end{theorem}

By the Jordan--von~Neumann--Wigner classification~\cite{JordanVonNeumannWigner1934}, the
simple Euclidean Jordan algebras are: $M_n(\mathbb{R})^{\mathrm{sa}}$,
$M_n(\mathbb{C})^{\mathrm{sa}}$, $M_n(\mathbb{H})^{\mathrm{sa}}$, the
spin factors $V_n$, and the exceptional Albert algebra
$M_3(\mathbb{O})^{\mathrm{sa}}$.

\begin{remark}[Non-associativity]
The sequential product $\sp{a}{b}$ is non-associative for general
(non-compatible) effects: there exist triples $(a,b,c)$ with
$\sp{(\sp{a}{b})}{c} \ne \sp{a}{(\sp{b}{c})}$. This is consistent with
S5, which requires associativity only for compatible pairs. By the theorem
of Westerbaan, Westerbaan, and van de
Wetering~\cite{WesterbaanWesterbaanVdW2020}, associativity would force
commutativity, which would force the state space to be a simplex
(classical). Non-associativity is thus a necessary feature of any
non-classical sequential product.
\end{remark}
